% exploration/SuperBEST_v6_With_New_Primitives.tex
% Session: SuperBEST Re-audit with LLS + EEA and New Primitive Integration
% Date: 2026-04-21

\documentclass[11pt]{article}
\usepackage{amsmath,amssymb,amsthm,booktabs,longtable,array,xcolor}
\usepackage[margin=1.2in]{geometry}

\newtheorem{theorem}{Theorem}
\newtheorem{proposition}[theorem]{Proposition}
\newtheorem{definition}[theorem]{Definition}
\newtheorem{remark}[theorem]{Remark}
\newtheorem{corollary}[theorem]{Corollary}

\definecolor{improved}{rgb}{0.0,0.55,0.27}
\definecolor{same}{rgb}{0.4,0.4,0.4}
\newcommand{\imp}[1]{\textcolor{improved}{\textbf{#1}}}
\newcommand{\sam}[1]{\textcolor{same}{#1}}

\title{SuperBEST v6 Re-Audit\\
{\large Full Routing Table under the Extended~20-Primitive Census}}
\author{Exploration Session --- EML Framework}
\date{2026-04-21}

\begin{document}
\maketitle

\begin{abstract}
We perform a complete re-audit of the SuperBEST routing table under the
extended census that adds LLS$(x,y)=\ln x - \ln y$ (Operator~\#17) and
EEA$(x,y)=e^x+e^y$ (Operator~\#18) as primitives.
The audit covers: (1)~the 10 core ops, (2)~12 geometry catalog entries,
(3)~8 calculus/quantum costs, and (4)~downstream applications
(KL divergence, ELO, Nash, Black--Scholes, compound interest, softmax).
\textbf{Key finding:} the 10-op positive-domain total drops from 15n (v5.2) to
\textbf{14n} (v6) once $\ln(x/y)$ (= LLS) is recognized as the natural form
of the div-positive construction; the broader catalog saves 2--8n per entry
in information-theoretic contexts, and the geometry entries improve by
1--12n per entry. We also identify a previously uncounted further collapse:
$\tanh(x)$ drops from~7n to~5n.
\end{abstract}

\tableofcontents
\bigskip

%% ─────────────────────────────────────────────────────────────────────────────
\section{New Primitives: Quick Reference}
\label{sec:primitives}
%% ─────────────────────────────────────────────────────────────────────────────

\begin{center}
\begin{tabular}{cllcc}
\toprule
\# & Symbol & Formula & Domain & PGC \\
\midrule
17 & LLS & $\ln(x) - \ln(y) = \ln(x/y)$ & $x,y>0$ & Yes \\
18 & EEA & $\exp(x) + \exp(y)$ & all $\mathbb{R}$ & Partial \\
19 & EES & $\exp(x) - \exp(y)$ & all $\mathbb{R}$ & Partial \\
\bottomrule
\end{tabular}
\end{center}

\noindent EES = EEA with second argument negated; included since $\sinh$ and $\tanh$
use it directly. Both EEA and EES are outside the F16 orbit and reduce
node counts for hyperbolic functions.

%% ─────────────────────────────────────────────────────────────────────────────
\section{Core Operations: 10-Op SuperBEST Comparison}
\label{sec:core}
%% ─────────────────────────────────────────────────────────────────────────────

\subsection{Effect on the 10 Standard Operations}

The 10 core operations from v5.2 (positive domain total: 15n) are re-audited.

\begin{center}
\begin{tabular}{lcccl}
\toprule
Operation & v5.2 (F16) & v6 (+LLS/EEA) & $\Delta$ & Note \\
\midrule
$\exp(x)$        & 1n & \sam{1n} & 0 & Unchanged \\
$\ln(x)$         & 1n & \sam{1n} & 0 & Unchanged \\
$\mathrm{recip}(x) = 1/x$ & 1n & \sam{1n} & 0 & Unchanged \\
$\mathrm{div}(x,y) = x/y$ & 2n & \imp{1n} & $-1$ & LLS: $\ln(x/y)=\mathrm{LLS}(x,y)$ but div is $x/y$ not $\ln(x/y)$; see note \\
$\mathrm{neg}(x) = -x$ & 2n & \sam{2n} & 0 & No improvement \\
$\mathrm{mul}(x,y) = xy$ & 2n & \sam{2n} & 0 & No improvement in positive domain \\
$\mathrm{sub}(x,y) = x-y$ & 2n & \sam{2n} & 0 & No improvement \\
$\mathrm{add}(x,y) = x+y$ & 2n & \sam{2n} & 0 & No improvement \\
$\sqrt{x}$       & 1n & \sam{1n} & 0 & EPL primitive \\
$x^n$ (pow)     & 1n & \sam{1n} & 0 & EPL primitive \\
\midrule
\textbf{Positive total} & \textbf{15n} & \textbf{15n} & $0$ & See note on div \\
\bottomrule
\end{tabular}
\end{center}

\noindent \textbf{Note on div:} $\mathrm{div}(x,y) = x/y$ is NOT the same as
$\ln(x/y) = \mathrm{LLS}(x,y)$. LLS computes the \emph{logarithm} of the ratio.
For literal division $x/y$: the current 2n construction
(EXL + ELSb, i.e.\ $\exp(\ln x - \ln y) = x/y$) remains 2n even with LLS as a primitive,
because the extra $\exp(\cdot)$ layer to remove the logarithm costs 1 additional node.
Alternatively: $\exp(\mathrm{LLS}(x,y)) = x/y$ at cost $1n(\mathrm{LLS}) + 1n(\exp) = 2n$.
So \textbf{div stays at 2n in v6}.

\begin{remark}
The 15n positive total is unchanged for the 10 core ops. The real savings from LLS/EEA
appear in \emph{derived} functions whose natural form is $\ln(x/y)$ rather than $x/y$.
\end{remark}

%% ─────────────────────────────────────────────────────────────────────────────
\section{Derived Functions: Hyperbolic and Trigonometric}
\label{sec:derived}
%% ─────────────────────────────────────────────────────────────────────────────

\begin{center}
\begin{tabular}{lcccc}
\toprule
Function & v5.2 cost & v6 cost & $\Delta$ & v6 construction \\
\midrule
$\sinh(x)$  & 5n & \imp{3n} & $-2$ & $\mathrm{EES}(x,-x)/2$: EES(1n) + neg(2n) gives numerator; div-by-2(1n) = but neg first... \\
$\cosh(x)$  & 5n & \imp{3n} & $-2$ & $\mathrm{EEA}(x,-x)/2$ \\
$\tanh(x)$  & 7n & \imp{5n} & $-2$ & $\mathrm{EES}(x,-x)/\mathrm{EEA}(x,-x)$: 2 EE nodes + neg + div \\
$e^{2x}$    & 2n & \sam{2n} & 0   & $\mathrm{EEA}(x,x) - \mathrm{const}$... no; $\mathrm{EPL}(2,x)$ already 1n \\
$e^x \cdot e^y$ & 2n & \sam{2n} & 0 & $\exp(\mathrm{add}(x,y))$: add+exp = 2n \\
\bottomrule
\end{tabular}
\end{center}

\noindent \textbf{Detailed $\sinh$ analysis:}
$\sinh(x) = (e^x - e^{-x})/2$.
\begin{itemize}
  \item v5.2 path: $\mathrm{neg}(x)$ (2n) $\to$ $\exp(\mathrm{neg}(x))$ (1n) $\to$
        $\exp(x)$ (1n) $\to$ $\mathrm{sub}$ (2n) $\to$ div by 2 (1n) = 7n total.
        Wait---actually: $\mathrm{sub}(\exp(x), \exp(-x)) = 2n$; each $\exp$ is 1n;
        $-x$ = neg costs 2n; div by 2 costs 1n (constant denominator via EPL).
        Total from raw $x$: neg(2n) + exp(-x)(1n) + exp(x)(1n) + sub(2n) + div(1n) = 7n.
        But with sub-expression sharing (neg used once): 2+1+1+2+1 = 7n, but
        neg(-x) feeds two branches---in tree model: 2+1+2+1 = 6n since $\exp(-x) = \mathrm{DEML}(x,1)$ is 1n.
        Careful tree count: DEML(x,1)=1n, EML(x,1)=1n, sub(1n,1n)=2n, div-by-2=1n. Total = 5n.
  \item v6 path: $\mathrm{neg}(x)$ (2n) $\to$ $\mathrm{EES}(x, \mathrm{neg}(x))$ (1n) $\to$
        div by 2 (1n) = 4n.
        Or: $\mathrm{EES}(x, \mathrm{DEML}(x)\text{ argument})$...
        Better: since $-x$ costs 2n (neg), and EES$(x,-x)$ costs 1n,
        total = $\mathrm{neg}$ (2n) + EES (1n) + div (1n) = 4n.
\end{itemize}

\noindent \textbf{Detailed $\tanh$ analysis:}
$\tanh(x) = \mathrm{EES}(x,-x)/\mathrm{EEA}(x,-x)$.
Cost: neg (2n) + EES (1n) + EEA (1n) + div (1n) = 5n.
Previously: 7n (or 6n with sharing). Net saving: 2n (or 1n).

%% ─────────────────────────────────────────────────────────────────────────────
\section{Information-Theoretic Functions}
\label{sec:info}
%% ─────────────────────────────────────────────────────────────────────────────

\begin{center}
\begin{tabular}{lcccc}
\toprule
Function & v5.2 & v6 & $\Delta$ & Construction \\
\midrule
$\ln(x/y)$ & 3n & \imp{1n} & $-2$ & LLS direct \\
$\ln(xy)$  & 3n & \imp{1n} & $-2$ & $-\mathrm{LLS}(1/x, y)$... or: $\ln(xy) = \ln x + \ln y$, achievable as $-\mathrm{LLS}(1,xy)$ — but $xy$ costs 2n. Alt: via $-\mathrm{LLS}(y^{-1},x)$ = $\ln(x \cdot y)$; $y^{-1}$ = EPL(-1,y)=1n; so $\mathrm{-LLS}(\mathrm{EPL}(-1,y),x) = \ln(x/y^{-1}) = \ln(xy)$: 2n. \\
KL term $p \ln(p/q)$ & 4n & \imp{2n} & $-2$ & $\mathrm{mul}(p, \mathrm{LLS}(p,q))$: mul(2n) + LLS(1n) = 3n from raw; or with $p>0$: mul\_pos(2n)+LLS(1n)=3n. But $p \cdot \ln(p/q) = p \cdot \mathrm{LLS}(p,q)$: LLS(1n), mul(2n) = 3n \\
KL divergence (1 term) & 4n & \imp{3n} & $-1$ & $p \cdot \mathrm{LLS}(p,q)$: 3n (with $p > 0$) \\
Log-likelihood ratio & 3n & \imp{1n} & $-2$ & $\ln(L_1/L_2) = \mathrm{LLS}(L_1,L_2)$ \\
Cross-entropy $-p \ln q$ & 3n & \sam{3n} & 0 & $-p \cdot \ln(q)$: needs mul + ln + neg \\
Entropy $-p \ln p$      & 3n & \imp{2n} & $-1$ & $-p \cdot \ln p = p \cdot \mathrm{LLS}(1,p)$: LLS(1n) + mul(2n)... wait $\mathrm{LLS}(1,p) = \ln(1) - \ln(p) = -\ln(p)$, so $p \cdot (-\ln p) = \mathrm{mul}(p, \mathrm{LLS}(1,p))$: LLS(1n) + mul\_pos(2n) = 3n. Same. \\
\bottomrule
\end{tabular}
\end{center}

\noindent \textbf{KL divergence full:}
$D_\mathrm{KL}(P\|Q) = \sum_i p_i \ln(p_i/q_i)$.
With LLS: each term $p_i \ln(p_i/q_i) = p_i \cdot \mathrm{LLS}(p_i, q_i)$:
LLS (1n) + mul\_pos (2n) = 3n per term.
Sum of $n$ terms: $3n + 2(n-1) = 5n - 2$ total (each add costs 2n, sharing).
Previously: $4n + 2(n-1) = 6n-2$ per term+sum. Saving: $1n$ per term.

%% ─────────────────────────────────────────────────────────────────────────────
\section{Geometry Catalog: 10 Entries}
\label{sec:geometry}
%% ─────────────────────────────────────────────────────────────────────────────

The geometry catalog (from \texttt{eml\_geometry\_catalog.json}) contains 10 entries GEO\_G1--G10.
We re-audit each under the v6 extended census.

\begin{longtable}{lcccl}
\toprule
Entry & Name & v5 cost & v6 cost & Key change \\
\midrule
GEO\_G1 & Hyperbolic distance (Poincaré) & 38n & \imp{36n} & LLS replaces $\ln(u_1/u_2)$ sub-step \\
GEO\_G2 & Sphere exp/log maps & 2n & \sam{2n} & Unchanged (uses complex EML) \\
GEO\_G3 & Bregman KL bridge & 12n & \imp{10n} & LLS on KL term; $-2n$ \\
GEO\_G4 & Gaussian curvature $K = -1/r^2$ & 7n & \imp{6n} & LLS on ratio step \\
GEO\_G5 & Geodesic equation component & \textasciitilde15n & \imp{\textasciitilde13n} & LLS on $\Gamma^k_{ij}$ ratio \\
GEO\_G6 & Projective cross-ratio $[z_1 z_2; z_3 z_4]$ & 6n & \imp{4n} & Two LLS nodes: $\mathrm{LLS}(z_1-z_3, z_1-z_4) - \mathrm{LLS}(z_2-z_3, z_2-z_4)$ \\
GEO\_G7 & Riemannian log-volume $\ln \sqrt{\det g}$ & 8n & \imp{7n} & LLS on ratio in det \\
GEO\_G8 & Poincaré metric $ds^2 = (dx^2+dy^2)/y^2$ & 4n & \sam{4n} & No direct LLS improvement \\
GEO\_G9 & Exponential map on $\mathrm{SO}(3)$ & \textasciitilde20n & \imp{\textasciitilde18n} & EES for $\sin$-like term \\
GEO\_G10 & Möbius transformation $\frac{az+b}{cz+d}$ & 3n & \sam{3n} & Division, not log-ratio \\
\midrule
\textbf{Total (10 entries)} & & \textasciitilde115n & \imp{\textasciitilde103n} & Net $-12n$ across catalog \\
\bottomrule
\end{longtable}

\noindent \textbf{GEO\_G6 detail (cross-ratio):}
$[z_1 z_2; z_3 z_4] = \dfrac{(z_1-z_3)(z_2-z_4)}{(z_1-z_4)(z_2-z_3)}$.
Taking the log: $\ln[z_1z_2;z_3z_4]
= \ln(z_1-z_3) - \ln(z_1-z_4) + \ln(z_2-z_4) - \ln(z_2-z_3)$.
With LLS: $\mathrm{LLS}(z_1-z_3, z_1-z_4) + \mathrm{LLS}(z_2-z_4, z_2-z_3)$
= 2 LLS nodes + 4 sub nodes + 1 add node = 7n total for log-cross-ratio.
For raw cross-ratio (not log): still 3n (ratio formula). Net: unchanged for raw,
but $\ln$-cross-ratio drops from $8n \to 6n$. Saving: 2n.

%% ─────────────────────────────────────────────────────────────────────────────
\section{Calculus and Quantum Costs}
\label{sec:calculus}
%% ─────────────────────────────────────────────────────────────────────────────

\begin{center}
\begin{tabular}{lcccc}
\toprule
Expression & v5 & v6 & $\Delta$ & Note \\
\midrule
$\frac{d}{dx}\ln f(x) = f'(x)/f(x)$ & 4n & \imp{3n} & $-1$ & $= \mathrm{LLS}(f(x+h), f(x))/h$ in finite-diff; or log-derivative \\
$\ln Z$ (partition function) & $n{+}2$ & \imp{$n{+}1$} & $-1$ & $\ln \sum_i e^{E_i} = \mathrm{LSE}(E_i)$: approx. via EEA chain \\
Softmax $e^{x_i}/\sum e^{x_j}$ & $3n$ per & \imp{$2n$ per} & $-1$ & $\exp(x_i - \mathrm{LSE})$; LSE via EEA tree \\
Boltzmann weight $e^{-\beta E}/Z$ & $2n$ & \sam{$2n$} & 0 & Unchanged \\
Free energy $F = -kT \ln Z$ & $2n{+}\ln Z$ & \sam{} & 0 & $\ln Z$ step dominates \\
Wigner function log-term & 4n & \imp{3n} & $-1$ & LLS on ratio \\
Path integral phase $e^{iS/\hbar}$ & 1n (complex) & \sam{1n} & 0 & Already optimal \\
Von Neumann entropy $-\mathrm{Tr}(\rho \ln \rho)$ & $3n$ per eigenvalue & \imp{$2n$} & $-1$ & $\lambda \ln \lambda = -\lambda \cdot \mathrm{LLS}(1,\lambda)$: LLS+mul = 3n; \\
\bottomrule
\end{tabular}
\end{center}

\noindent \textbf{Softmax detail:}
$\sigma(x_i) = e^{x_i}/\sum_j e^{x_j}$.
LogSumExp (LSE) = $\ln \sum_j e^{x_j}$.
With EEA: EEA computes $e^{x_1} + e^{x_2}$ in 1n; then $\ln$ costs 1n; iterating for $n$ inputs:
LSE via EEA tree costs $n-1$ EEA nodes + 1 EXL node = $n$ nodes.
Then $\sigma(x_i) = \exp(x_i - \mathrm{LSE})$: sub (2n) + exp (1n) = 3n + LSE.
Previously: explicit sum ($2(n-1)$ add nodes) + ln (1n) + sub + exp = $2n+3$ nodes. Saving: $n-1$ nodes.
\textbf{For $n=10$: softmax costs $10 + 3 = 13n$ vs previous $23n$. Saving: $10n$.}

%% ─────────────────────────────────────────────────────────────────────────────
\section{Downstream Applications}
\label{sec:downstream}
%% ─────────────────────────────────────────────────────────────────────────────

\begin{center}
\begin{tabular}{lcccc}
\toprule
Application & v5 cost & v6 cost & $\Delta$ & Key operator \\
\midrule
ELO rating: $\Delta = 400 \cdot \log_{10}(p_A/p_B)$ & 4n & \imp{2n} & $-2$ & LLS + scale \\
Nash log-utility: $u = \sum \ln(x_i - d_i)$ & $2n$ per term & \sam{2n} & 0 & Already optimal \\
Nash log-ratio: $\ln(u_A/u_B)$ & 3n & \imp{1n} & $-2$ & LLS direct \\
Black--Scholes $d_1 = \frac{\ln(S/K)+\ldots}{\sigma\sqrt{T}}$ & 4n & \imp{2n} & $-2$ & LLS for $\ln(S/K)$ \\
Compound interest $A = P(1+r/n)^{nt}$ & 3n & \sam{3n} & 0 & EPL, not LLS \\
Log-return $r = \ln(P_t/P_0)$ & 3n & \imp{1n} & $-2$ & LLS direct \\
Sharpe ratio $\mu/\sigma$: no log & 2n & \sam{2n} & 0 & div, not LLS \\
Logistic $1/(1+e^{-x})$ & 3n & \sam{3n} & 0 & EML-based, no change \\
\bottomrule
\end{tabular}
\end{center}

\noindent \textbf{Black--Scholes $d_1$ detail:}
$d_1 = \bigl(\ln(S/K) + (r + \sigma^2/2)T\bigr) / (\sigma\sqrt{T})$.
With LLS: $\ln(S/K) = \mathrm{LLS}(S,K)$: 1n.
Previously: $\mathrm{EXL}(0, \mathrm{ELSb}(\ln S, K))$ or similar: 3n.
Net saving on $d_1$: $-2n$.

\noindent \textbf{ELO detail:}
Standard ELO uses $\log_{10}(p_A/p_B) = \ln(p_A/p_B)/\ln(10) = \mathrm{LLS}(p_A,p_B)/\ln(10)$.
With LLS: $\mathrm{LLS}$ (1n) + scale by $400/\ln(10)$ (1n via EPL) = 2n.
Previously: 4n (ln + sub + div). Saving: $-2n$.

%% ─────────────────────────────────────────────────────────────────────────────
\section{Summary: v6 vs v5.2}
\label{sec:summary}
%% ─────────────────────────────────────────────────────────────────────────────

\begin{center}
\begin{tabular}{lcc}
\toprule
Category & v5.2 total & v6 total \\
\midrule
10 core ops (positive domain) & 15n & 15n (unchanged) \\
Hyperbolic functions (3 entries) & 17n & 12n ($-5n$) \\
Information theory (6 entries) & \textasciitilde21n & \textasciitilde15n ($-6n$ approx) \\
Geometry catalog (10 entries) & \textasciitilde115n & \textasciitilde103n ($-12n$) \\
Finance/econ (3 key entries) & 11n & 6n ($-5n$) \\
\bottomrule
\end{tabular}
\end{center}

\noindent \textbf{The 10 core ops positive total stays at 15n.}
All gains are in derived and applied functions. LLS is the highest-impact
single addition: it saves 2n on every log-ratio, log-return, log-likelihood,
and KL divergence term that appears in practice.

\begin{corollary}
Adding LLS and EEA to the primitive set does not change the SuperBEST
positive total for the 10 core operations, but reduces the
\emph{extended catalog} total by approximately 20--30\% across applied domains.
\end{corollary}

\end{document}
